\section{Experimental Design}
\label{sec:exp}

Because recursive retry already achieves high success on split-knowledge tasks, we treat task success as a \emph{constraint} rather than the primary optimization objective. Our evaluation is therefore a controlled ablation: we isolate the effect of Bayesian belief-guided routing within an otherwise identical recursive delegation loop, measuring whether informed agent selection improves token usage, agent calls, and time-to-first-success relative to uninformed (uniform random) selection.

\begin{figure*}[t]
\centering
\resizebox{\linewidth}{!}{%
\begin{tikzpicture}[
  >=latex, thick,
  box/.style={rectangle, draw, rounded corners, align=center, font=\small,
              minimum width=28mm, minimum height=8mm, fill=blue!6},
  metric/.style={rectangle, draw, rounded corners, align=left, font=\scriptsize,
                 minimum width=36mm, minimum height=6mm, fill=green!5},
  line/.style={-latex}
]

% RQs
\node[box] (rq1) at (0,0) {RQ1: Composition Efficiency};
\node[box] (rq2) at (0,-2.2) {RQ2: Specialization};
\node[box] (rq3) at (0,-4.4) {RQ3: Adaptability};

% Hypotheses
\node[box, right=3.8cm of rq1] (h1) {H1: Belief-guided routing reduces\\ cost at matched success};
\node[box, right=3.8cm of rq2] (h2) {H2: Belief scores diverge,\\ reflecting specialization};
\node[box, right=3.8cm of rq3] (h3) {H3: Routing adapts to impaired\\ agents in real time};

% Metrics
\node[metric, right=5.2cm of h1] (m1) {
\begin{tabular}{@{}l@{}}
-- Success Rate (constraint) \\
-- Tokens \\
-- Agent Calls \\
-- Time-to-first-success
\end{tabular}};
\node[metric, right=5.2cm of h2] (m2) {
\begin{tabular}{@{}l@{}}
-- Belief Score Evolution \\
-- Rounds to Expert Selection
\end{tabular}};
\node[metric, right=5.2cm of h3] (m3) {
\begin{tabular}{@{}l@{}}
-- Impairment Test (pre/post) \\
-- Belief Score Decline \\
-- Contribution Frequency \\
-- Final Output Score
\end{tabular}};

% Connections
\draw[line] (rq1) -- (h1);
\draw[line] (rq2) -- (h2);
\draw[line] (rq3) -- (h3);

\draw[line] (h1) -- (m1);
\draw[line] (h2) -- (m2);
\draw[line] (h3) -- (m3);

\end{tikzpicture}%
}
\caption{\textbf{Evaluation framework.} Research questions, testable hypotheses, and associated metrics. Each dimension targets a distinct property of the routing mechanism.}
\label{fig:rq-flow}
\end{figure*}

\subsection{Research Questions}
We organize evaluation around three dimensions (Figure~\ref{fig:rq-flow}):
\begin{itemize}[leftmargin=1.4em]
    \item[\textbf{RQ1}] \textbf{Composition efficiency (H1):} At matched task success, does belief-guided routing reduce token usage, agent calls, and time-to-first-success compared to random recursive delegation?
    \item[\textbf{RQ2}] \textbf{Specialization (H2):} Do belief scores diverge over task sequences, with domain experts being preferentially selected for related tasks?
    \item[\textbf{RQ3}] \textbf{Adaptability (H3):} When a previously reliable agent is impaired mid-experiment, does the routing policy detect the degradation and redirect queries?
\end{itemize}

\subsection{Tasks and Agent Population}
We constructed a benchmark of split-knowledge tasks via a three-stage pipeline (full details in Appendix): (i)~domain-specific question generation for each agent (20--22 questions each), (ii)~multi-agent task synthesis by sampling 4--6 agents and composing questions that require non-redundant contributions from each, and (iii)~iterative sampling until the question pool was exhausted. By construction, each task requires contributions from multiple agents and cannot be solved by any single agent alone.

The agent population comprises 6 specialist RAG agents (biology, finance, law, medicine, electrical engineering, mathematics) with retrieval over curated knowledge bases, and 10 generalist conversational agents (fitness, literature, technology, geography, storytelling, politics, academics, career guidance, trivia, daily planning). All agents use Azure OpenAI endpoints with fixed random seeds for reproducibility. Agents are initialized with uniform priors $\alpha_0=\beta_0=1$ unless memory-aware initialization is enabled.

\subsection{Baseline}
Since split-knowledge tasks are unsolvable by any single agent, single-agent baselines are excluded by task design. The comparison instead isolates the routing mechanism:
\begin{itemize}[leftmargin=1.4em]
    \item \textbf{Random Delegation:} The full recursive loop (judging, re-routing, aggregation, budget limits) is retained identically, but agent selection at each step is uniformly random. This ensures that any observed difference is attributable to belief-guided routing rather than to recursion, judging, or aggregation.
\end{itemize}

\subsection{Evaluation Metrics}
\textbf{Performance (constraint):} Task success rate (final score $\geq 85$ as judged by a calibrated LLM) and output quality, reported to confirm that both methods achieve comparable success.
\textbf{Efficiency (H1):} Token usage, agent calls, and time-to-first-success, reported as ratios normalized to Random Delegation (lower is better).
\textbf{Specialization (H2):} Belief score trajectories over task sequences and rounds-to-expert-selection for domain-specific tasks.
\textbf{Adaptability (H3):} In an impairment test (50 normal followed by 50 degraded tasks for one agent), we track belief score decline, contribution frequency shift, and system output quality.
Full metric definitions are provided in the Appendix.
